% Part: methods
% Chapter: induction
% Section: strong-induction

\documentclass[../../../include/open-logic-section]{subfiles}

\begin{document}

\olfileid{mth}{ind}{str}

\olsection{قوي استقرا}

د استقرا په هغه اصل کښې چې پورته مو بيان کړ، $P(0)$
ثابتوو، او دا هم چې که $P(k)$ وي نو $P(k+1)$ هم وي۔ په
دويمه برخه کښې فرض کوو چې $P(k)$ رښتيا دے، او د دې
فرض په کارولو $P(k+1)$ ثابتوو۔ البته، په معادل ډول
$P(k-1)$ هم فرضولاے شو او د $P(k)$ د ثابتولو دپاره
ترې کار اخيستلاے شو۔ مهمه خبره دا ده چې له هر عدد
څخه ئې راتلونکي ته استدلال کولاے شو؛ يعنې د هر عدد
دپاره اړونده دعوه په دې فرض ثابتولاے شو چې د هغه د
مخکيني عدد دپاره سمه ده۔

د استقرا د اصل يوه بله بڼه شته چې په هغې کښې يوازې
دا نۀ فرضوو چې دعوه د مخکيني $k-1$ دپاره سمه
ده، چې د $k$ مخکينے عدد دے، بلکې د~$k$ څخه د ټولو کوچنيو عددونو دپاره ئې
فرضوو، او د دې فرض په کارولو دعوه د~$k$ دپاره
ثابتوو۔ دا هم دعوه $P(n)$ راکوي، د ټولو~$n \in
\Nat$ دپاره۔ ځکه چې يو ځل $P(0)$
ثابت کړو، نو په دې سره مو ثابته کړې وي چې $P$ له~$1$
څخه د ټولو کوچنيو عددونو دپاره سم دے۔ او که پوهېږو
چې له $P(l)$ څخه، د ټولو $l<k$ دپاره، $P(k)$ راوځي، نو
په ځانګړي ډول د $k=1$ دپاره هم په دې پوهېږو۔ نو
$P(1)$ نتيجه اخيستلاے شو۔ په دې سره مو $P(0)$ او
$P(1)$ ثابت کړي، يعنې $P(l)$، د ټولو $l<2$ دپاره۔ او
ځکه چې دا شرط هم لرو چې که $P(l)$ وي، د ټولو $l<2$ دپاره،
نو $P(2)$ هم وي، نو~$P(2)$ نتيجه اخيستلاے شو، او
همداسې مخ په وړاندې ځو۔

په حقيقت کښې، که دا عام شرط ثابتولاے شو چې «د ټولو
$k$ دپاره، که $P(l)$ وي، د ټولو $l<k$ دپاره، نو $P(k)$
هم وي»، بيا د $P(0)$ جلا ثابتولو ته اړتيا نشته، ځکه
چې له همدغه شرطه راوځي۔ په ياد ولرئ چې داسې عامه
دعوه لکه «د ټولو $l<k$ دپاره $P(l)$» هغه وخت رښتيا
ده چې هېڅ $l<k$ نۀ وي۔ دا د تشې کميت‌ښودنې يو حالت
دے: «ټول $A$، $B$ دي» هغه وخت رښتيا ده چې هېڅ $A$
نۀ وي؛ $\lforall[x][(!A(x) \lif !B(x))]$ هغه وخت
رښتيا دے چې هېڅ $x$،~$!A(x)$ نۀ پوره کوي۔ په دې
حالت کښې ئې رسمي بڼه «$\lforall[l][(l < k \lif P(l))]$»
ده، او دا هغه وخت رښتيا ده چې هېڅ $l < k$ نۀ وي۔
او که $k=0$ وي، همدا حالت دے: هېڅ $l<0$ نشته؛ نو
«د ټولو $l<0$ دپاره $P(0)$» رښتيا دے، که $P$ هر
څه وي۔ له دې امله د دې شرط ثبوت چې «که $P(l)$ وي، د ټولو
$l<k$ دپاره، نو $P(k)$ هم وي»، په خپله~$P(0)$
ثابتوي۔

دا بڼه هغه وخت ګټوره ده چې د~$k$ دپاره د دعوې
ثابتول يوازې د $k-1$ په دعوه تکيه نۀ شي کولاے،
بلکې ښايي دې فرض ته اړتيا ولري چې د يوه يا څو
$l<k$ دپاره رښتيا ده۔

\end{document}
